\documentclass[10pt, conference, letterpaper]{IEEEtran}

% ==========================================
% PACKAGES
% ==========================================
\usepackage[utf8]{inputenc}
\usepackage[T1]{fontenc}
\usepackage{graphicx}
\usepackage{amsmath, amssymb}
\usepackage{booktabs}
\usepackage{cite}
\usepackage{hyperref}

% ==========================================
% FRONT MATTER
% ==========================================
\title{From Reactive to Proactive: An AI-Driven Early Warning System for Enhancing Supply Chain Resilience Through News-Operational Data Fusion}

\author{
\IEEEauthorblockN{Duc-Duong Dao Minh\IEEEauthorrefmark{1}, Buu-Tanh Tran-Le, Huyen-Huynh Chau Nhu, Thuy-Nguyen Nhut, Linh-Le Quynh Khanh}
\IEEEauthorblockA{\textit{University of Economics and Law} \\
\textit{Ho Chi Minh City, Vietnam} \\
\IEEEauthorrefmark{1}Corresponding author: ducdmk24406@st.uel.edu.vn}
}

\begin{document}

\maketitle

\begin{abstract}
% [PASTE YOUR ABSTRACT TEXT HERE]
% Ensure the abstract is concise and summarizes the core problem, methodology, and findings.
Existing supply chain disruption forecasting approaches remain confined to a single data modality, leaving a critical blind-spot: no validated pipeline demonstrates whether fusing heterogeneous modalities yields measurable Incremental Predictive Validity over unimodal baselines. This study proposes a four-stage Cascading AI architecture that transforms raw logistics news into calibrated risk features via DistilBERT-based binary filtering and zero-shot multi-label classification, then fuses them with structured ERP data through a Geographic Weighting function before training gradient-boosted classifiers. Experiments on 8,728 logistics news articles (2022--2024) and five aerospace operational datasets yield three principal findings: (1) the Gatekeeper binary filter achieves ROC-AUC = 0.8927 and Recall = 0.9503; (2) the fused XGBoost model attains Precision = 0.1654 on minority-class stockout prediction --- a 28.7\% relative improvement over the ERP-only baseline, substantiating Incremental Predictive Validity; and (3) the system provides a Lead-Time Gain of 1.0 to 3.2 weeks across all component families prior to disruption onset. These results validate that integrating public-domain textual signals with operational telemetry produces an actionable early warning capability for upstream inventory risk.
\end{abstract}

\begin{IEEEkeywords}
% Keywords are separated by em-dashes (---) as per IEEE conventions.
Supply Chain Risk Management---Natural Language Processing---Heterogeneous Data Fusion---Early Warning System---Inventory Disruption Forecasting
\end{IEEEkeywords}

% ==========================================
% BODY
% ==========================================

\section{Introduction}
% [PASTE INTRODUCTION TEXT HERE]

\subsection{Motivation}
% [PASTE MOTIVATION TEXT HERE]

\subsection{Research Problem}
% [PASTE RESEARCH PROBLEM TEXT HERE]

\subsection{Research Gap}
% [PASTE RESEARCH GAP TEXT HERE]

% PLACEHOLDER FOR TABLE 1: Positioning Matrix comparing our study with Cano-Marin (2023) and Ivanov (2022)
\begin{table*}[!t]
\centering
\caption{Positioning Matrix comparing our study with Cano-Marin (2023) and Ivanov (2022)}
\label{tab:positioning}
\begin{tabular}{@{}llll@{}}
\toprule
\textbf{Feature/Dimension} & \textbf{Cano-Marin (2023)} & \textbf{Ivanov (2022)} & \textbf{Proposed Study} \\ \midrule
% [PASTE TABLE CONTENT HERE]
\bottomrule
\end{tabular}
\end{table*}

\section{Literature Review}
% [PASTE LITERATURE REVIEW TEXT HERE]

\subsection{NLP in SCRM}
% [PASTE TEXT HERE]

\subsection{Machine Learning Pitfalls}
% [PASTE TEXT HERE: Specifically mentioning Temporal Leakage and Spurious Regression]

\section{Methodology}
% [PASTE METHODOLOGY INTRO TEXT HERE]

\subsection{Phase 0: Sensing Layer}
% [PASTE PHASE 0 TEXT HERE]

\subsection{Phase 1: Sense-making Layer}
% [PASTE PHASE 1 TEXT HERE]

\subsection{Phase 2: Feature Fusion and Machine Learning}
% [PASTE PHASE 2 TEXT HERE]

\subsection{Phase 3: Evaluation and Data Governance}
% [PASTE PHASE 3 TEXT HERE: Mention SHA-256 Checksumming for data integrity]

% PLACEHOLDER FOR FIGURE 1: Research Methodology Framework (Rigor-focused)
\begin{figure*}[!t]
\centering
\includegraphics[width=\textwidth]{diagram/figure_1_methodology.png}
\caption{Research Methodology Framework (Rigor-focused).}
\label{fig:methodology}
\end{figure*}

% PLACEHOLDER FOR FIGURE 2: Cascading System Architecture (3-layer technical flow)
\begin{figure*}[!t]
\centering
\includegraphics[width=\textwidth]{diagram/figure_2_architecture.png}
\caption{Cascading System Architecture (3-layer technical flow).}
\label{fig:architecture}
\end{figure*}

\section{Results \& Discussion}
% [PASTE RESULTS TEXT HERE]
% INCLUDE: Strategic Calibration of PR-AUC vs ROC-AUC

% PLACEHOLDER FOR 3-Tier Evaluation Matrix (Rule-Based vs. Baseline ML vs. SCRM ML)
\begin{table*}[!t]
\centering
\caption{3-Tier Evaluation Matrix}
\label{tab:evaluation}
\begin{tabular}{@{}llll@{}}
\toprule
\textbf{Metric} & \textbf{Tier 1 (Rule-Based)} & \textbf{Tier 2 (Baseline ML)} & \textbf{Tier 3 (SCRM ML)} \\ \midrule
% [PASTE TABLE CONTENT HERE]
\bottomrule
\end{tabular}
\end{table*}

\section{Conclusion}
% [PASTE CONCLUSION TEXT HERE]
% Summarizing the Actionable early warning window and Lead-time Gains

% ==========================================
% REFERENCES
% ==========================================
\bibliographystyle{IEEEtran}
\bibliography{references} % Create a references.bib file for your citations

\end{document}
